\documentclass[11pt]{article}
\usepackage{amsmath,amssymb,amsthm}
\usepackage{geometry}
\usepackage{enumitem}
\usepackage{setspace}

\geometry{margin=1in}

\title{CSE400: Fundamentals of Probability in Computing\\
Lecture 6: Discrete Random Variables, Expectation, and Problem Solving}
\author{Dhaval Patel, PhD}
\date{January 22, 2025}

\begin{document}

\maketitle

\hrule
\vspace{0.5em}

\section{Random Variables: Motivation and Concept}

\subsection*{Definition (Slide)}

A \textbf{random variable} \( X \) on a sample space \( \Omega \) is a function
\[
X : \Omega \rightarrow \mathbb{R}
\]
that assigns to each sample point \( \omega \in \Omega \) a real number \( X(\omega) \).

\subsection*{Clarification (Handwritten Notes)}

\begin{itemize}
    \item A \emph{random variable is not a variable} in the algebraic sense; \textbf{it is a function}.
    \item The randomness lies in the outcome \( \omega \), not in the function itself.
\end{itemize}

\subsection*{Discrete Restriction (Slide)}

In this lecture, attention is restricted to \textbf{discrete random variables}, i.e.,
\begin{itemize}
    \item Although \( X \) maps into \( \mathbb{R} \),
    \item The actual set of values \( \{ X(\omega) : \omega \in \Omega \} \) is \textbf{finite or countably infinite}.
\end{itemize}


\vspace{1em}

\section{Visualizing a Discrete Random Variable}

\subsection*{Distribution as a Bar Diagram (Slide)}

\begin{itemize}
    \item \textbf{x-axis}: values that the random variable can take
    \item \textbf{Height of bar at value \( a \)}:
    \[
    \Pr[X = a]
    \]
    \item Each probability is computed from the probability of the \textbf{corresponding event} in the sample space.
\end{itemize}

\vspace{1em}

\section{Types of Random Variables}

\subsection{Discrete Random Variables (Slides + Notes)}

A \textbf{discrete random variable} has:
\begin{itemize}
    \item \textbf{Countable support}
    \item \textbf{Probability Mass Function (PMF)}
    \item Probabilities assigned to \textbf{single values}
    \item Each possible value has \textbf{strictly positive probability}
\end{itemize}

\subsubsection*{Common Discrete Random Variables (Listed)}

\begin{itemize}
    \item Bernoulli
    \item Binomial
    \item Geometric
    \item Poisson
    \item Hypergeometric
\end{itemize}


\subsection{Continuous Random Variables (Slides + Notes)}

A \textbf{continuous random variable} has:
\begin{itemize}
    \item \textbf{Uncountable support}
    \item \textbf{Probability Density Function (PDF)}
    \item Probabilities assigned to \textbf{intervals}
    \item Each individual value has \textbf{zero probability}
\end{itemize}

\vspace{1em}

\section{Example 1: Tossing 3 Fair Coins}

\subsection*{Experiment Description (Slide)}

\begin{itemize}
    \item Toss \textbf{3 fair coins}
    \item Let \( Y \) = number of heads observed
\end{itemize}

\subsection*{Possible Values}

\[
Y \in \{0,1,2,3\}
\]

\subsection*{Probability Computation (Slide)}

\[
\Pr(Y=0)=\Pr(t,t,t)=\frac{1}{8}
\]
\[
\Pr(Y=1)=\Pr(t,t,h),\Pr(t,h,t),\Pr(h,t,t)=\frac{3}{8}
\]
\[
\Pr(Y=2)=\Pr(h,h,t),\Pr(h,t,h),\Pr(t,h,h)=\frac{3}{8}
\]
\[
\Pr(Y=3)=\Pr(h,h,h)=\frac{1}{8}
\]

\subsection*{PMF Summary}

\[
\Pr(Y=0)=\frac{1}{8},\;
\Pr(Y=1)=\frac{3}{8},\;
\Pr(Y=2)=\frac{3}{8},\;
\Pr(Y=3)=\frac{1}{8}
\]

\subsection*{Legitimacy Condition (Slide)}

Since \( Y \) must take one of these values:
\[
\sum_{i=0}^{3} \Pr(Y=i)=1
\]

\vspace{1em}

\section{Probability Mass Function (PMF)}

\subsection*{Definition (Slide)}

A random variable that can take on \textbf{at most a countable number of values} is said to be \textbf{discrete}.

Let \( X \) be a discrete random variable with range
\[
R_X = x_1, x_2, x_3, \ldots
\quad \text{(finite or countably infinite)}
\]

The function
\[
P_X(x_k) = \Pr(X = x_k), \quad k=1,2,3,\ldots
\]
is called the \textbf{Probability Mass Function (PMF)} of \( X \).

\subsection*{Normalization Condition (Slide + Notes)}

Since \( X \) must take one of the values \( x_k \),
\[
\sum_{k=1}^{\infty} P_X(x_k) = 1
\]

This expresses that the sample space is \textbf{collectively exhaustive}.

\vspace{1em}

\section{PMF Example with Parameter \( \lambda \)}

\subsection*{Given (Slide)}

The PMF of a random variable \( X \) is:
\[
p(i) = c \frac{\lambda^i}{i!}, \quad i = 0,1,2,\ldots
\]
where \( \lambda > 0 \).

Find:
\begin{enumerate}
    \item \( \Pr(X=0) \)
    \item \( \Pr(X>2) \)
\end{enumerate}

\subsection*{Step 1: Find \( c \)}

Since PMF must sum to 1:
\[
\sum_{i=0}^{\infty} p(i) = 1
\]
\[
c \sum_{i=0}^{\infty} \frac{\lambda^i}{i!} = 1
\]

Using:
\[
e^{\lambda} = \sum_{i=0}^{\infty} \frac{\lambda^i}{i!}
\]

Thus:
\[
c e^{\lambda} = 1
\quad \Rightarrow \quad
c = e^{-\lambda}
\]

\subsection*{Step 2: Compute \( \Pr(X=0) \)}

\[
\Pr(X=0) = e^{-\lambda} \frac{\lambda^0}{0!} = e^{-\lambda}
\]

\subsection*{Step 3: Compute \( \Pr(X>2) \)}

\[
\Pr(X>2) = 1 - \Pr(X \le 2)
\]
\[
= 1 - [\Pr(X=0) + \Pr(X=1) + \Pr(X=2)]
\]
\[
= 1 - \left(
e^{-\lambda}
+ e^{-\lambda}\lambda
+ e^{-\lambda}\frac{\lambda^2}{2}
\right)
\]

\vspace{1em}

\section{Bayes’ Theorem (Recap)}

\subsection*{Bayes Formula (Slide)}

Using:
\[
\Pr(A \cap B_i) = \Pr(B_i|A)\Pr(A)
\]

We obtain:
\[
\Pr(B_i|A)
=
\frac{\Pr(A|B_i)\Pr(B_i)}
{\sum_{j=1}^{n} \Pr(A|B_j)\Pr(B_j)}
\]

This is referred to as \textbf{Bayes Formula (Proposition 3.1)}.

\subsection*{Terminology (Slide)}

\begin{itemize}
    \item \( \Pr(B_i) \): \textbf{Prior probability}
    \item \( \Pr(B_i|A) \): \textbf{Posterior probability}
\end{itemize}

\vspace{1em}

\section{Bayes’ Theorem Example: Auditorium with 30 Rows}

\subsection*{Setup (Slide)}

\begin{itemize}
    \item Auditorium has \textbf{30 rows}
    \item Row 1 has 11 seats
    \item Row 2 has 12 seats
    \item Row 3 has 13 seats
    \item \(\ldots\)
    \item Row 30 has 40 seats
\end{itemize}

\subsection*{Prize Selection Process}

\begin{enumerate}
    \item A row is selected \textbf{uniformly at random}
    \[
    \Pr(R_k) = \frac{1}{30}
    \]
    \item A seat is selected \textbf{uniformly at random within that row}
\end{enumerate}

\subsection*{Question 1}

Compute:
\[
\Pr(S_{15} \mid R_{20})
\]

\subsection*{Solution (Slide + Notes)}

\begin{itemize}
    \item Row 20 has \textbf{30 seats}
    \item Only one seat is seat 15
\end{itemize}

\[
\Pr(S_{15} \mid R_{20}) = \frac{1}{30}
\]

\subsection*{Question 2}

Compute:
\[
\Pr(R_{20} \mid S_{15})
\]

\subsection*{Using Bayes’ Formula}

\[
\Pr(R_{20} \mid S_{15})
=
\frac{\Pr(S_{15} \mid R_{20}) \Pr(R_{20})}
{\Pr(S_{15})}
\]

Where:
\[
\Pr(S_{15}) = \sum_{k=15}^{30} \Pr(S_{15} \mid R_k)\Pr(R_k)
\]


\hrule
\vspace{1em}

\section*{End of Lecture 6}

\end{document}
